\documentclass[english, 12pt, a4paper, elec, utf8, a-2b, online]{aaltothesis}

\setupthesisfonts

\usepackage{iftex}
\ifXeTeX
  \usepackage{fontspec}
  \usepackage{xeCJK}
  \setCJKmainfont[AutoFakeBold=2,AutoFakeSlant=.2]{FandolSong-Regular}
  \setCJKsansfont[AutoFakeBold=2,AutoFakeSlant=.2]{FandolHei-Regular}
  \setCJKmonofont[AutoFakeBold=2,AutoFakeSlant=.2]{FandolFang-Regular}
\else
  \usepackage{CJKutf8}
\fi
\usepackage{graphicx}
\graphicspath{{images/}{../images/}{latex/../images/}}
\usepackage{array}
\usepackage{booktabs}
\usepackage{siunitx}
\usepackage[type={CC}, modifier={by-nc-sa}, version={4.0}]{doclicense}
\usepackage[style=ieee,sorting=none]{biblatex}
\IfFileExists{thesisreferences.bib}{%
  \addbibresource{thesisreferences.bib}%
}{%
  \IfFileExists{../thesisreferences.bib}{%
    \addbibresource{../thesisreferences.bib}%
  }{%
    \addbibresource{latex/../thesisreferences.bib}%
  }%
}

\urlstyle{tt}
\DeclareRobustCommand{\codepath}[1]{\nolinkurl{#1}}
\pdfstringdefDisableCommands{\def\codepath#1{#1}}
\newcolumntype{L}[1]{>{\raggedright\arraybackslash}p{#1}}
\setlength{\headheight}{36pt}
\setlength{\emergencystretch}{2em}
\newcommand{\thesisinput}[1]{%
  \IfFileExists{#1.tex}{%
    \input{#1}%
  }{%
    \IfFileExists{latex/#1.tex}{%
      \input{latex/#1}%
    }{%
      \input{../latex/#1}%
    }%
  }%
}

\degreeprogram{Electronics and Nanotechnology}
\major{Microelectronic Circuit Design}
\univdegree{MSc}
\thesisauthor{Linan Jiang}
\thesistitle{基于 Mie 散射成像的数据驱动穿深预测与概率性碰壁筛查}
\thesissubtitle{面向多孔柴油喷雾的不确定性感知代理模型工作流}
\place{Espoo}
\date{19 May 2026}
\supervisor{Prof.\ Ossi Kaario}
\advisor{Dr Zeeshan Ahmad}
\advisor{Dr Qiang Cheng}
\collaborativepartner{Wärtsilä}
\uselogo{'}
\copyrighttext{\noexpand\textcopyright\ \number\year. This work is
    licensed under a Creative Commons "Attribution-NonCommercial-ShareAlike 4.0
    International" (BY-NC-SA 4.0) license.}{\noindent\textcopyright\ \number
    \year\ \doclicenseThis}
\keywords{Mie 散射\spc 喷雾成像\spc 穿深\spc 计算机视觉\spc 代理模型}
\thesisabstract{多燃料船用发动机中的柴油预喷具有持续时间短、瞬态特征强等特点。若液相包络在形成有效点火核之前即到达活塞顶或缸壁，可能引发壁面润湿、润滑油稀释以及混合准备恶化等问题。本文建立了一套基于实验证据的筛查工作流，将高速 Mie 散射视频转化为以穿深为核心的几何观测量，并构建用于早期喷雾撞壁风险筛查的代理模型。研究数据来自 Wärtsilä 光学喷雾燃烧舱的室温、非反应、顶视 Mie 记录，覆盖多种工况与喷嘴几何。该工作流结合人工喷嘴标定、CUDA 加速的多孔视频处理、逐喷油束特征提取，以及一种用于平滑噪声标签并缓解视场边界右删失低估偏差的降阶轨迹拟合方法。基于拟合目标，本文训练了三阶段神经网络代理模型：阶段一学习确定性代表轨迹基线；阶段二扩展至完整过滤数据集，并预测异方差均值与方差；阶段三在早期观测可靠时对齐原始 CDF 数据，在晚期观测稀疏时保留保守的教师引导，以避免将离开视场误判为物理饱和。在来自 9 个喷嘴试验批次的 542,565 个预测点上（未删失 CDF 评估协议），最终无锚点阶段三模型（均值 MSE 与对数方差 MSE 的混合蒸馏，方差权重 5.0）在五种随机种子自举下的 RMSE 为 4.27 mm（标准差 0.04 mm），MAE 为 3.03 mm（标准差 0.04 mm），并保持 88.7 percent / 98.9 percent 的 1-sigma / 2-sigma 不确定性覆盖率，相对校准后的 Hiroyasu--Arai 和 Naber--Siebers 物理基线 RMSE 降低超过 54 percent。本文同时在相同协议下评估了稀疏异方差 Gaussian process（SVGP）作为替代主干：分布内与留一喷嘴族协议上 SVGP 是更强的逐点回归器（RMSE 4.19 mm 与跨族 5.95 mm 对 MLP 的 7.00 mm），而本文 MLP 的结构性留任理由是更保守的不确定性覆盖、用于碰壁筛查接口的喷射起始辅助输出，以及消融实验在硕士级算力预算下的可负担性；二者应被理解为同一删失感知学习协议下不同的精度-覆盖率权衡，而非主干层面的一边倒优势。本文进一步通过简化几何投影，将预测穿深分布映射为概率性撞壁筛查代理指标。需要强调的是，该工作流的适用范围限于测量域内的早期工程筛查，不能替代 CFD、发动机试验或经验证的缸内撞壁物理模型。}

\begin{document}
\ifXeTeX
\else
\begin{CJK*}{UTF8}{gbsn}
\fi

\makecoverpage
\makecopyrightpage

\begin{abstractpage}[english]
多燃料船用发动机中的柴油预喷具有持续时间短、瞬态特征强等特点。若液相包络在形成有效点火核之前即到达活塞顶或缸壁，可能引发壁面润湿、润滑油稀释以及混合准备恶化等问题。本文建立了一套基于实验证据的筛查工作流，将高速 Mie 散射视频转化为以穿深为核心的几何观测量，并构建用于早期喷雾撞壁风险筛查的代理模型。

研究数据来自 Wärtsilä 光学喷雾燃烧舱（OSCC）的室温、非反应、顶视 Mie 记录，覆盖多种工况与喷嘴几何。该工作流结合人工喷嘴标定、CUDA 加速的多孔视频处理、逐喷油束特征提取，以及一种用于平滑噪声标签并缓解视场边界右删失低估偏差的降阶轨迹拟合方法。

基于拟合目标，本文训练了三阶段神经网络代理模型。Stage~1 学习确定性代表轨迹基线。Stage~2 扩展至完整过滤数据集，并预测异方差均值与方差。Stage~3 在早期观测可靠时对齐原始 CDF 数据，在晚期观测稀疏时保留保守的教师引导，以避免将离开视场误判为物理饱和。

在来自 9 个喷嘴试验批次的 542{,}565 个预测点上（未删失 CDF 评估协议），最终无锚点阶段三模型（均值 MSE 与对数方差 MSE 的混合蒸馏，方差权重 $\lambda_\sigma{=}5.0$）在五种随机种子自举下的 RMSE 为~\SI{4.27}{mm}（标准差~\SI{0.04}{mm}），MAE 为~\SI{3.03}{mm}（标准差~\SI{0.04}{mm}），并保持 $88.7\,\%$ / $98.9\,\%$ 的保守 1$\sigma$/2$\sigma$ 不确定性覆盖率，相对校准后的 Hiroyasu--Arai 和 Naber--Siebers 物理基线 RMSE 降低超过 54\%。本文同时评估了稀疏异方差 Gaussian process（SVGP）作为替代主干：在分布内的未删失协议上 SVGP 点精度略优（RMSE~\SI{4.19}{mm}），在留一喷嘴族协议下点精度优势进一步扩大（剔除脱离设计族的 Nozzle0 后，SVGP 四折均值 RMSE~\SI{5.95}{mm} 对 MLP 的~\SI{7.00}{mm}）。本文 MLP 的结构性留任理由是更保守的不确定性覆盖、用于碰壁筛查接口的喷射起始辅助输出，以及消融实验在硕士级算力预算下的可负担性；二者应被理解为同一删失感知学习协议下不同的精度--覆盖率权衡，而非主干层面的一边倒优势。本文进一步通过简化几何投影，将预测穿深分布映射为概率性撞壁筛查代理指标。需要强调的是，该工作流的适用范围限于测量域内的早期工程筛查，不能替代 CFD、发动机试验或经验证的缸内撞壁物理模型。
\end{abstractpage}

\mysection{前言}
谨向 Ossi Kaario 教授致以诚挚谢意，感谢其在本论文研究过程中的指导。感谢 Zeeshan Ahmad 博士与 Qiang Cheng 博士在喷雾成像、数据处理和建模方面给予的技术支持。亦感谢参与喷雾处理流程协作的各位同事，他们的讨论与反馈帮助我不断明确本研究的工程问题与实现路径。

\mysection{AI 工具使用声明}
本论文在准备过程中使用了 AI 辅助工具，用于探索性代码实现、脚本自动化、数据筛查辅助、代码与处理流程总结、可视化准备、辅助作图、英文改写、语法修正以及语言润色。AI 工具被纳入分层且可审计的研究工作流中使用：代码层面的输出通过可执行程序、日志和中间文件验证；数据与图表层面的输出与导出表格、诊断图和数值结果核对；文本层面的改写则以前述代码、数据、图表、文献和作者解释为依据。

AI 工具未被作为科学依据、引用来源、实验结果来源或独立作者使用。本文的研究问题、方法选择、实现决策、结果解释及最终文本均由作者审阅、修改并确认。作者对本文的准确性、原创性、代码、分析、图表、结果与结论承担全部责任。

\thesistableofcontents
\dothesispagenumbering{}

\thesisinput{sections_zh/01_introduction}
\thesisinput{sections_zh/02_literature_review}
\thesisinput{sections_zh/03_data_workflow}
\thesisinput{sections_zh/04_trajectory_surrogate_screening}
\thesisinput{sections_zh/05_results}
\thesisinput{sections_zh/06_discussion}
\thesisinput{sections_zh/07_conclusion}
\thesisinput{sections_zh/appendix_implementation_mapping}
\thesisbibliography
\printbibliography

\ifXeTeX
\else
\end{CJK*}
\fi
\end{document}
